%!TEX root = ../dokumentation.tex

\addchap{\langabkverz}
%nur verwendete Akronyme werden letztlich im Abkürzungsverzeichnis des Dokuments angezeigt
%Verwendung: 
%		\ac{Abk.}    --> fügt die Abkürzung ein, beim ersten Aufruf wird zusätzlich automatisch die ausgeschriebene Version davor eingefügt bzw. in einer Fußnote (hierfür muss in header.tex \usepackage[printonlyused,footnote]{acronym} stehen) dargestellt
%		\acs{Abk.}   -->  fügt die Abkürzung ein
%		\acf{Abk.}    --> fügt die Abkürzung UND die Erklärung ein
%		\acl{Abk.}   --> fügt nur die Erklärung ein
%		\acp{Abk.}  --> gibt Plural aus (angefügtes 's'); das zusätzliche 'p' funktioniert auch bei obigen Befehlen
%	siehe auch: http://golatex.de/wiki/%5Cacronym
%	
\begin{acronym}[seccomp]
    \setlength{\itemsep}{0pt}
    \acro{ad}[AD]{Active-Directory}
    \acro{aes}[AES]{Advanced Encryption Standard}
    \acro{app}[App]{Application}
    \acro{ard}[ARD]{Arbeitsgemeinschaft der öffentlich-rechtlichen Rundfunkanstalten der Bundesrepublik Deutschland}
    \acro{aws}[AWS]{Amazon Web Services}
    
    \acro{bpf}[BPF]{Berkeley Packet Filter}
    \acro{br}[BR]{Bayerischer Rundfunk}
    
    \acro{cbo}[CBO]{Cloud Break-Out}
    \acro{cgroups}[Cgroups]{Control Groups}
    \acro{cli}[CLI]{Command Line Interface}
    \acro{cpu}[CPU]{Central processing unit}
    \acro{cvss}[CVSS]{Common Vulnerability Scoring System}
    
    \acro{das}[DAS]{Direct-Attach-Storage}
    \acro{dr}[DR]{Deutschlandradio}
    \acro{dw}[DW]{Deutsche Welle}
    
    \acro{ft}[FT]{Fachteam}
    
    \acro{gb}[GB]{Gigabyte}
    \acro{gcp}[GCP]{Google Cloud Platform}
    \acro{ghz}[GHz]{Gigahertz}
    \acro{gui}[GUI]{Graphical User Interface}
    
    \acro{ha}[HA]{Hauptabteilung}
    \acro{hp}[HP]{Hewlett-Packard}
    \acro{hr}[HR]{Hessischer Rundfunk}
    \acro{hsb}[HSB]{Hauptstadtstudio Berlin}
    \acro{http}[HTTP]{Hypertext Transfer Protocol}
    \acro{https}[HTTPS]{Hypertext Transfer Protocol Secure}
    
    \acro{iaas}[IaaS]{Infrastructure as a Service}
    \acro{ieee}[IEEE]{Institute of Electrical and Electronics Engineers}
    \acro{imp}[IMP]{IT, Medien- und Produktionstechnik}
    \acro{ip}[IP]{Internet Protocol}
    \acro{ipsec}[IPsec]{Internet Protocol Security}
    \acro{iot}[IoT]{Internet of Things}
    \acro{isp}[ISP]{Internet Service Provider}
    \acro{it}[IT]{Informationstechnik}
    
    \acro{json}[JSON]{JavaScript Object Notation}
    
    \acro{kb}[KB]{Kilobyte}
    
    \acro{lan}[LAN]{Local Area Network}
    \acro{lbo}[LBO]{Local Internet Breakout}
    \acro{lte}[LTE]{Long Term Evolution}
    
    \acro{mac}[MAC]{Media Access Control}
    \acro{macc}[MAC]{Mandatory Access Control}
    \acro{mb}[MB]{Megabyte}
    \acro{mdm}[MDM]{Mobile Device Management}
    
    \acro{nas}[NAS]{Network attached Storage}
    \acro{ndr}[NDR]{Norddeutscher Rundfunk}
    \acro{nist}[NIST]{National Institute of Standards and Technology}
    \acro{nvme}[NVMe]{Non Volatile Memory Express}
    
    \acro{os}[OS]{Operating System}
    
    \acro{paas}[PaaS]{Platform as a Service}
    \acro{pc}[PC]{Personal Computer}
    \acro{pdf}[PDF]{Portable Document Format}
    
    \acro{qos}[QoS]{Quality of Service}
    
    \acro{rb}[RB]{Radio Bremen}
    \acro{rfa}[RFA]{Rundfunkanstalten}
    \acro{rj}[RJ]{Registered Jack}
    
    \acro{s}[s]{Sekunde}
    \acro{saas}[SaaS]{Software as a Service}
    \acro{selinux}[SELinux]{Security-Enhanced Linux}
    \acro{seccomp}[seccomp]{Secure Computing}
    \acro{sd-wan}[SD-WAN]{Software-Defined Wide Area Network}
    \acro{sdn}[SDN]{Software Definied Network}
    \acro{sgx}[SGX]{Software Guard Extensions}
    \acro{sla}[SLA]{Service Level Agreement}
    \acro{ssh}[SSH]{Secure Shell}
    \acro{ssd}[SSD]{Solid State Drive}
    \acro{sso}[SSO]{Single Sign-On}
    \acro{sr}[SR]{Saarländischer Rundfunk}
    \acro{stp}[STP]{Sterpunkt}
    \acro{suid}[SUID]{Set User ID}
    \acro{swr}[SWR]{Südwestrundfunk}
    
    \acro{tb}[TB]{Terabyte}
    \acro{tcp}[TCP]{Transmission Control Protocol}
    
    \acro{ui}[UI]{User Interface}
    \acro{uid}[UID]{User Identification}
    
    \acro{vlan}[VLAN]{Virtual Local Area Network}
    \acro{voip}[VoIP]{Voice over Internet Protocol}
    \acro{vpn}[VPN]{Virtual Private Network}
    \acro{vm}[VM]{Virtuelle Maschine}
    \acroplural{vm}[VMs]{Virtuelle Maschinen}
    
    \acro{wan}[WAN]{Wide Area Network}
    \acro{wdr}[WDR]{Westdeutscher Rundfunk}
    \acro{wlan}[WLAN]{Wireless Local Area Network}

    \acro{yaml}[YAML]{yet another markup language}
    
    \acro{zbs}[ZBS]{Zentraler Beitragsservice}
\end{acronym}
